\section{数据设计}

\subsection{数据设计目标}
本系统的数据设计需要同时满足当前图像检测业务的稳定运行和后续论文、评审文本检测能力的扩展需要。由于现有实现已经围绕图像上传、检测任务、检测结果、人工审核、通知和日志形成了可运行的数据链路，本轮概要设计不直接推翻既有模型，而是在保留现有实体的基础上抽象统一资源、统一任务和统一结果结构。

数据设计的主要目标如下：
\begin{itemize}[leftmargin=2em]
    \item 支撑用户、组织、资源、检测任务、检测结果、人工审核、报告、通知和日志等核心业务对象的持久化；
    \item 保持对现有图像检测表结构的兼容，避免在迭代初期破坏已有上传、检测、审核和报告流程；
    \item 为论文检测和评审检测新增数据对象预留扩展空间，使三类学术资源能够纳入统一任务和统一报告链路；
    \item 保证资源归属、组织边界、任务状态、结果数据和日志记录之间具备可追溯关系；
    \item 支持异步检测、批量检测、人工审核、报告生成和管理端统计等场景下的数据查询与状态流转。
\end{itemize}

\subsection{数据模型设计原则}
本系统数据模型遵循以下原则：
\begin{itemize}[leftmargin=2em]
    \item \textbf{统一抽象与差异化载荷分离}：用户、组织、资源、任务、结果等通用数据应统一建模；图像、论文和评审文本的专有结果通过差异化结果对象或扩展字段保存；
    \item \textbf{现有结构优先复用}：当前已经存在的 \texttt{User}、\texttt{Organization}、\texttt{FileManagement}、\texttt{ImageUpload}、\texttt{DetectionTask}、\texttt{DetectionResult} 等实体应继续作为近期实现基础；
    \item \textbf{任务链路可追溯}：从资源上传到检测任务、检测结果、人工审核、报告文件和操作日志之间必须保留可回查的关联关系；
    \item \textbf{状态字段集中管理}：资源、检测任务、检测结果、审核任务和通知均应有明确状态枚举，避免前后端对同一状态产生不同解释；
    \item \textbf{敏感资源权限隔离}：论文正文、评审文本、图像文件和检测报告均属于学术资源，应通过用户归属和组织归属限制访问范围。
\end{itemize}

\subsection{核心数据模型类图}
\begin{figure}[H]
    \centering
    \resizebox{0.98\textwidth}{!}{\begin{tikzpicture}[
  font=\scriptsize,
  every node/.style={align=center},
  class/.style={
    rectangle split,
    rectangle split parts=3,
    draw,
    rounded corners=2pt,
    minimum width=3.45cm,
    text width=3.05cm,
    inner xsep=4pt,
    inner ysep=3pt,
    fill=blue!5
  },
  rel/.style={-{Latex[length=2mm]}, thick},
  evt/.style={-{Latex[length=2mm]}, thick, dashed},
  legend/.style={draw, rounded corners=2pt, fill=gray!8, align=left, text width=5.2cm, inner sep=5pt}
]

\node[class, fill=blue!8] (org) at (0,0) {组织\nodepart{second}
组织名称\\管理员\\剩余配额
\nodepart{third}
重置配额\\查询额度};

\node[class, fill=green!8] (user) at (4.8,0) {用户\nodepart{second}
用户名\\角色\\权限位\\组织ID
\nodepart{third}
鉴权\\校验权限};

\node[class, fill=yellow!10] (file) at (9.6,0) {文件资源\nodepart{second}
文件名\\文件类型\\标签
\nodepart{third}
登记文件\\查询文件};

\node[class, fill=yellow!10] (task) at (0,-3.9) {检测任务\nodepart{second}
任务名\\状态\\检测参数
\nodepart{third}
创建任务\\更新状态};

\node[class, fill=orange!10] (image) at (4.8,-3.9) {图像资源\nodepart{second}
来源文件\\页码\\检测标记
\nodepart{third}
登记图像\\标记状态};

\node[class, fill=orange!12] (result) at (9.6,-3.9) {检测结果\nodepart{second}
真假判定\\置信度\\检测时间
\nodepart{third}
保存结果\\汇总结果};

\node[class, fill=orange!12] (sub) at (12.2,-7.8) {子检测结果\nodepart{second}
方法名\\概率\\掩码数据
\nodepart{third}
记录子方法};

\node[class, fill=red!10] (request) at (0,-7.8) {审核请求\nodepart{second}
申请时间\\发布者状态\\管理员状态
\nodepart{third}
提交申请\\分派审核人};

\node[class, fill=red!10] (manual) at (4.8,-7.8) {人工审核\nodepart{second}
审核人\\审核状态\\报告文件
\nodepart{third}
提交审核\\生成报告};

\node[class, fill=red!10] (imgreview) at (9.6,-7.8) {图像审核项\nodepart{second}
分值\\理由\\标注点\\结论
\nodepart{third}
记录人工评分};

\node[class, fill=purple!10] (log) at (4.8,-11.6) {日志与通知\nodepart{second}
操作类型\\关联对象\\消息状态
\nodepart{third}
记录日志\\发送通知};

\node[legend] (legend) at (10.6,-11.6) {图例：\\
实线表示主业务数据关联；\\
虚线表示日志或通知等事件记录关系。};

\draw[rel] (org) -- (user);
\draw[rel] (user) -- (file);
\draw[rel] (file) -- (image);
\draw[rel] (user) -- (task);
\draw[rel] (task) -- (image);
\draw[rel] (image) -- (result);
\draw[rel] (result) -- (sub);
\draw[rel] (result) -- (request);
\draw[rel] (request) -- (manual);
\draw[rel] (manual) -- (imgreview);
\draw[evt] (task) -- (log);
\draw[evt] (request) -- (log);
\draw[evt] (manual) -- (log);

\end{tikzpicture}
}
    \caption{核心数据模型 UML 类图}
\end{figure}

图 5-1 展示了系统核心数据对象之间的静态关系。图中用户、组织、文件、图像、检测任务、检测结果、人工审核、日志和通知对应当前代码中已经存在的主要模型；统一资源、论文资源、评审资源、论文检测结果、评审检测结果和综合报告属于后续迭代中建议新增或抽象出的对象。该图的目的不是给出最终数据库物理表，而是明确数据层演进方向：在保留图像检测闭环的同时，将论文与 Review 纳入统一资源和统一任务体系。

\subsection{核心实体说明}
\begin{longtable}{|p{4.8cm}|p{2cm}|p{7.3cm}|}
\hline
实体名称 & 类型 & 说明 \\ \hline
\texttt{User} & 已有实体 & 系统用户主体，记录账号、邮箱、角色、头像、组织归属、权限位、个人检测次数等信息，是资源、任务、审核、日志和通知的主要归属对象。 \\ \hline
\texttt{Organization} & 已有实体 & 组织主体，记录组织名称、邮箱、管理员、证明材料、组织级检测配额等信息，用于支持组织内资源管理和权限隔离。 \\ \hline
\texttt{OrganizationApplication} & 已有实体 & 组织入驻申请记录，保存申请状态、组织信息、管理员账号信息、证明材料和审核人。 \\ \hline
\texttt{InvitationCode} & 已有实体 & 组织邀请码记录，保存邀请码、所属组织、角色、有效期和使用状态。 \\ \hline
\texttt{FileManagement} & 已有实体 & 上传文件元数据记录，保存文件名称、大小、类型、标签、上传用户和组织归属。 \\ \hline
\texttt{ImageUpload} & 已有实体 & 图像资源记录，保存图像文件、来源文件、是否从 PDF 提取、页码、检测状态、审核状态和图像级判定。 \\ \hline
\texttt{DetectionTask} & 已有实体 & 检测任务主体，保存任务名称、发起用户、所属组织、状态、上传时间、完成时间、报告文件和检测参数。 \\ \hline
\texttt{DetectionResult} & 已有实体 & 图像检测结果主体，保存是否疑似造假、置信度、检测时间、状态、LLM 判断、ELA 图像、LLM 可视化图像和 EXIF 判定等信息。 \\ \hline
\texttt{SubDetectionResult} & 已有实体 & 图像子检测结果，保存某一子检测方法的概率、掩膜图像和掩膜矩阵。 \\ \hline
\texttt{ReviewRequest} & 已有实体 & 人工审核申请，保存申请用户、关联检测结果、待审图像集合、申请原因、审核起止时间、管理员审批状态和审核人员集合。 \\ \hline
\texttt{ManualReview} & 已有实体 & 人工审核执行记录，保存审核人员、审核状态、关联图像、审核报告文件和审核时间。 \\ \hline
\texttt{ImageReview} & 已有实体 & 图像人工审核明细，保存每种检测方法对应的人工评分、理由、标注点集和最终人工判断。 \\ \hline
\texttt{Feedback} & 已有实体 & 人工审核反馈记录，保存评论、点赞状态、反馈用户和反馈时间。 \\ \hline
\texttt{Log} & 已有实体 & 操作日志记录，保存用户、操作类型、关联模型、关联记录 ID 和操作时间。 \\ \hline
\texttt{Notification} & 已有实体 & 通知消息记录，保存发送方、接收方、通知类型、标题、内容、状态、跳转链接和通知时间。 \\ \hline
\texttt{AcademicResource} & 预计新增 & 统一学术资源抽象，建议用于统一表示图像、论文和 Review 三类资源，并保存资源类型、原始文件、解析状态和归属信息。 \\ \hline
\texttt{PaperResource} & 预计新增 & 论文资源扩展对象，建议保存论文标题、正文文本、段落索引、解析状态和来源文件。 \\ \hline
\texttt{ReviewResource} & 预计新增 & 评审文本资源扩展对象，建议保存 Review 文本内容、语言类型、来源文件、解析状态和可选论文关联信息。 \\ \hline
\texttt{PaperDetectionResult} & 预计新增 & 论文检测结果扩展对象，建议保存全文级风险、AI 贡献占比、段落级概率、事实性鉴伪子结论和原因说明。 \\ \hline
\texttt{ReviewDetectionResult} & 预计新增 & Review 检测结果扩展对象，建议保存 AI 生成风险、模板化倾向、可疑语句、相似文本提示和置信度。 \\ \hline
\texttt{Report} & 抽象 & 报告元数据对象，当前可复用任务或人工审核中的报告文件字段；后续可独立抽象为检测报告、人工审核报告和综合报告的统一记录。 \\ \hline
\end{longtable}

\subsection{核心实体关系}
系统核心实体之间的关系如下：
\begin{itemize}[leftmargin=2em]
    \item 一个 \texttt{Organization} 可关联多个 \texttt{User}，一个用户在当前设计中最多归属于一个组织；
    \item 一个 \texttt{OrganizationApplication} 由管理员审核，审核通过后可生成对应的 \texttt{Organization} 和组织管理员用户关系；
    \item 一个 \texttt{InvitationCode} 关联一个组织和一个目标角色，用于将用户加入指定组织；
    \item 一个 \texttt{User} 可上传多个 \texttt{FileManagement} 文件记录，文件同时可带有组织归属；
    \item 一个 \texttt{FileManagement} 可派生多个 \texttt{ImageUpload}，例如从 PDF 或压缩包中提取多张图像；
    \item 一个 \texttt{DetectionTask} 可关联多个 \texttt{ImageUpload} 和多个 \texttt{DetectionResult}，用于支持批量图像检测；
    \item 一个 \texttt{DetectionResult} 关联一个图像检测对象，并可拥有多个 \texttt{SubDetectionResult} 子方法检测结果；
    \item 一个 \texttt{ReviewRequest} 可关联一个自动检测结果和多张待审核图像，并可分派给多个审核人员；
    \item 一个 \texttt{ManualReview} 关联一个 \texttt{ReviewRequest} 和一个审核人员，并可包含多个 \texttt{ImageReview} 明细；
    \item \texttt{Log} 通过关联模型名和关联记录 ID 记录上传、检测、审核申请和人工审核等业务事件；
    \item \texttt{Notification} 面向用户记录站内通知和状态提醒，可通过 URL 字段指向任务、审核或报告详情；
    \item 后续新增的 \texttt{AcademicResource} 应作为统一资源入口，与 \texttt{DetectionTask} 建立关联；图像、论文、Review 的专属资源对象作为其子类型或扩展对象存在；
    \item 后续论文和 Review 的检测结果建议以 \texttt{DetectionResult} 的统一结果主体为入口，再分别关联 \texttt{PaperDetectionResult} 和 \texttt{ReviewDetectionResult} 等差异化结果载荷。
\end{itemize}

\subsection{主要数据项设计}
概要设计阶段不展开全部数据库字段，但需要明确各类数据对象应包含的关键数据项。

\subsubsection{用户与组织数据}
\begin{longtable}{|p{4.8cm}|p{9.1cm}|}
\hline
数据对象 & 关键数据项 \\ \hline
\texttt{User} & 用户名、邮箱、密码摘要、角色、权限位、组织 ID、头像、简介、个人 LLM 与非 LLM 剩余次数、密码重置验证码和过期时间。 \\ \hline
\texttt{Organization} & 组织名称、邮箱、管理员用户 ID、创建时间、组织描述、Logo、证明材料、组织 LLM 与非 LLM 剩余次数、配额重置时间。 \\ \hline
\texttt{OrganizationApplication} & 申请组织名称、申请邮箱、管理员用户名、管理员邮箱、证明材料、Logo、描述、提交时间、申请状态、审核人。 \\ \hline
\texttt{InvitationCode} & 邀请码、组织 ID、目标角色、是否已使用、过期时间。 \\ \hline
\end{longtable}

\subsubsection{资源数据}
\begin{longtable}{|p{4.8cm}|p{9.1cm}|}
\hline
数据对象 & 关键数据项 \\ \hline
\texttt{FileManagement} & 上传用户、所属组织、文件名、文件大小、文件类型、上传时间、学科或资源标签。 \\ \hline
\texttt{ImageUpload} & 任务 ID、文件记录 ID、图像文件路径、是否从 PDF 提取、来源页码、上传时间、是否已检测、是否已提交审核、AI 检测是否判定为可疑。 \\ \hline
\texttt{AcademicResource} & 资源 ID、资源类型、原始文件 ID、所属用户、所属组织、解析状态、存储路径、创建时间、更新时间、是否归档。 \\ \hline
\texttt{PaperResource} & 资源 ID、论文标题、正文文本、段落索引、解析状态、来源文件 ID、解析错误信息。 \\ \hline
\texttt{ReviewResource} & 资源 ID、Review 文本、语言类型、来源文件 ID、解析状态、可选论文关联信息。 \\ \hline
\end{longtable}

\subsubsection{任务与结果数据}
\begin{longtable}{|p{4.8cm}|p{9.1cm}|}
\hline
数据对象 & 关键数据项 \\ \hline
\texttt{DetectionTask} & 任务名称、发起用户、所属组织、任务状态、上传时间、完成时间、报告文件、检测参数、是否启用 LLM。 \\ \hline
\texttt{DetectionResult} & 图像 ID、任务 ID、是否可疑、置信度、检测时间、检测状态、是否进入人工审核、LLM 判断、ELA 图像、LLM 图像、EXIF 判断。 \\ \hline
\texttt{SubDetectionResult} & 检测结果 ID、子方法编号、概率、掩膜图、掩膜矩阵、创建时间。 \\ \hline
\texttt{PaperDetectionResult} & 任务 ID、论文资源 ID、全文风险等级、AI 贡献占比、结果摘要、段落级概率列表、事实性鉴伪子结论、原因说明、模型版本。 \\ \hline
\texttt{ReviewDetectionResult} & 任务 ID、Review 资源 ID、AI 生成风险、模板化风险、可疑句段、相似文本提示、置信度、模型版本。 \\ \hline
\end{longtable}

\subsubsection{审核、报告、日志与通知数据}
\begin{longtable}{|p{4.8cm}|p{9.1cm}|}
\hline
数据对象 & 关键数据项 \\ \hline
\texttt{ReviewRequest} & 自动检测结果 ID、申请用户、所属组织、待审图像集合、申请时间、申请原因、发布者状态、管理员审批状态、审核起止时间、审核人员集合、管理员审核理由。 \\ \hline
\texttt{ManualReview} & 审核申请 ID、审核人员、所属组织、审核状态、审核图像集合、图像审核明细集合、审核时间、审核报告文件。 \\ \hline
\texttt{ImageReview} & 人工审核 ID、图像 ID、子方法评分、子方法理由、人工标注点集、最终人工判断、审核时间。 \\ \hline
\texttt{Feedback} & 人工审核 ID、反馈用户、是否点赞、评论内容、反馈时间。 \\ \hline
\texttt{Report} & 报告 ID、任务 ID、报告类型、报告文件路径、生成时间、生成状态、报告版本、所属用户和组织。 \\ \hline
\texttt{Log} & 操作时间、操作用户、操作类型、关联模型名称、关联记录 ID。 \\ \hline
\texttt{Notification} & 接收人、发送人、通知类别、标题、内容、阅读状态、通知时间、跳转链接。 \\ \hline
\end{longtable}

\subsection{状态设计}
状态字段是异步任务和人工审核流程的核心。系统应统一维护状态枚举，并在前后端共享状态含义。

\subsubsection{任务状态}
\begin{longtable}{|p{3cm}|p{10.8cm}|}
\hline
状态 & 含义 \\ \hline
\texttt{pending} & 任务已创建但尚未进入检测执行阶段。 \\ \hline
\texttt{in\_progress} & 任务正在执行检测、解析、结果生成或报告生成。 \\ \hline
\texttt{completed} & 任务已完成，检测结果和报告数据已可读取。 \\ \hline
\texttt{failed} & 建议新增状态，表示任务执行失败，需记录失败原因以供用户提示和管理员排查。 \\ \hline
\texttt{cancelled} & 建议新增状态，表示任务被用户或管理员取消。 \\ \hline
\end{longtable}

\subsubsection{检测结果状态}
\begin{longtable}{|p{3cm}|p{10.8cm}|}
\hline
状态 & 含义 \\ \hline
\texttt{in\_progress} & 检测结果正在生成，前端可展示处理中状态。 \\ \hline
\texttt{completed} & 检测结果已生成，用户可查看详情并下载报告。 \\ \hline
\texttt{failed} & 建议新增状态，表示检测模型调用、文件解析或结果写入失败。 \\ \hline
\end{longtable}

\subsubsection{人工审核状态}
\begin{longtable}{|p{3cm}|p{10.8cm}|}
\hline
状态 & 含义 \\ \hline
\texttt{pending} & 审核申请已提交，等待管理员审批或等待审核人员处理。 \\ \hline
\texttt{accepted} & 管理员已接受人工审核申请。 \\ \hline
\texttt{refused} & 管理员拒绝人工审核申请。 \\ \hline
\texttt{in\_progress} & 审核任务已分派给审核人员，正在处理。 \\ \hline
\texttt{completed} & 审核任务已完成，人工审核结论和报告已生成。 \\ \hline
\texttt{undo} & 审核人员尚未完成被分配的审核任务。 \\ \hline
\end{longtable}

\subsubsection{通知状态}
\begin{longtable}{|p{3cm}|p{10.8cm}|}
\hline
状态 & 含义 \\ \hline
\texttt{unread} & 通知已生成但用户尚未读取。 \\ \hline
\texttt{read} & 用户已读取通知。 \\ \hline
\end{longtable}

\subsection{数据流转设计}
系统主要数据流转过程如下：
\begin{enumerate}[leftmargin=2em]
    \item 用户通过用户端上传学术资源，系统创建文件管理记录，并根据资源类型生成图像资源记录或后续统一资源记录；
    \item 用户发起检测后，系统创建 \texttt{DetectionTask}，记录发起人、组织、任务参数和初始状态；
    \item 异步任务读取资源数据并调用图像、论文或 Review 检测模块，执行完成后写入检测结果或对应扩展结果对象；
    \item 对于图像检测，系统继续写入 \texttt{SubDetectionResult}，保存子方法概率、掩膜图和掩膜矩阵；
    \item 检测任务完成后，报告服务读取任务和结果数据生成报告文件，并将报告路径写回任务或独立报告对象；
    \item 用户申请人工审核时，系统创建人工审核申请，管理员审批后生成审核执行记录和审核明细；
    \item 关键操作写入 \texttt{Log}，任务状态变化或审核结果生成时写入 \texttt{Notification}；
    \item 用户端或管理端通过任务 ID、资源 ID、审核 ID 或日志条件查询对应数据，完成结果展示、统计和追溯。
\end{enumerate}

\subsection{兼容与演进设计}
当前系统已形成以图像为中心的数据结构，因此数据演进应分阶段完成：
\begin{itemize}[leftmargin=2em]
    \item \textbf{第一阶段：复用现有图像链路}。继续使用 \texttt{FileManagement}、\texttt{ImageUpload}、\texttt{DetectionTask}、\texttt{DetectionResult} 和 \texttt{SubDetectionResult} 支撑图像检测、人工审核和报告生成；
    \item \textbf{第二阶段：引入资源类型字段或统一资源表}。在不破坏现有图像流程的前提下，增加资源类型、解析状态和原始文件关联，使论文和 Review 能进入统一资源登记流程；
    \item \textbf{第三阶段：扩展结果载荷}。将现有图像结果保留为图像专用结果，同时新增论文检测结果和 Review 检测结果对象，并通过统一任务 ID 进行聚合；
    \item \textbf{第四阶段：统一报告和日志}。将检测报告、人工审核报告和综合鉴伪报告抽象为统一报告记录，同时扩展日志类型，支持论文检测日志和 Review 检测日志管理；
    \item \textbf{第五阶段：治理能力统一}。管理端资源检索、任务统计、模型配置和日志查询均基于资源类型与任务类型筛选，不再只围绕图像数据组织。
\end{itemize}

\subsection{数据完整性与约束}
为保证系统数据一致性，应在详细设计和实现阶段重点考虑以下约束：
\begin{itemize}[leftmargin=2em]
    \item 用户邮箱、组织名称、组织邮箱、邀请码等字段应保持唯一性；
    \item 任务、资源、结果和报告之间应有明确外键或逻辑关联，避免报告文件无法追溯到原始任务；
    \item 子检测结果应保证同一检测结果下同一检测方法只保存一条有效记录；
    \item 人工审核结果必须关联到明确的审核申请和审核人员，人工评分、理由和标注点集应与图像或文本片段保持对应关系；
    \item 删除资源、任务或组织时应采用限制删除、级联删除或软删除策略，并在详细设计中明确适用场景；
    \item 异步任务失败时应保证状态最终可落库，避免任务长期停留在处理中。
\end{itemize}

\subsection{索引与查询设计}
系统常见查询场景包括用户个人任务列表、组织任务统计、管理端日志检索、审核任务列表和检测结果详情查询。建议在以下字段上设置索引或优化查询：
\begin{itemize}[leftmargin=2em]
    \item 用户、组织外键字段，例如资源、任务、审核和日志中的用户 ID 与组织 ID；
    \item 状态字段，例如任务状态、审核状态、通知状态；
    \item 时间字段，例如上传时间、检测完成时间、审核时间、日志时间和通知时间；
    \item 资源类型和任务类型字段，用于支持图像、论文、Review 分类检索；
    \item 日志关联字段，例如关联模型名称和关联记录 ID；
    \item 报告文件关联字段，用于从任务详情快速定位报告。
\end{itemize}

\subsection{数据安全与权限控制}
论文、图像、Review 和检测报告均属于敏感学术资源，数据层应配合业务层实现以下安全要求：
\begin{itemize}[leftmargin=2em]
    \item 所有资源、任务、结果和报告都应记录所属用户和所属组织，查询时必须校验访问者是否具备访问权限；
    \item 普通用户只能访问本人资源、本人任务、被分派的审核任务及权限范围内的公开信息；
    \item 管理员操作组织、用户、审核和日志数据时应写入操作日志；
    \item 论文正文、Review 文本和报告文件不应在日志中直接明文记录，日志只记录引用 ID 和操作类型；
    \item 文件存储路径应由后端控制，不允许前端直接拼接物理文件路径访问资源；
    \item 重要删除操作建议采用软删除或归档策略，保留必要审计线索。
\end{itemize}

\subsection{本章小结}
本章从概要设计角度给出了系统数据模型的总体方案。当前实现中已有的数据模型可以支撑用户、组织、图像资源、检测任务、图像检测结果、人工审核、通知和日志等能力；后续迭代需要在此基础上引入统一资源抽象和面向论文、Review 的专用结果载荷。该设计能够在不破坏现有图像检测闭环的前提下，将系统逐步扩展为支持图像、论文和同行评审 Review 的综合学术 AI 鉴伪平台。
